\documentclass[11pt]{article}

\usepackage[a4paper,margin=1in]{geometry}
\usepackage{amsmath,amssymb,amsthm,mathtools}
\usepackage{graphicx}
\usepackage{hyperref}

\newtheorem{theorem}{Theorem}[section]
\newtheorem{proposition}[theorem]{Proposition}
\newtheorem{lemma}[theorem]{Lemma}
\newtheorem{corollary}[theorem]{Corollary}
\theoremstyle{definition}
\newtheorem{definition}[theorem]{Definition}
\newtheorem{remark}[theorem]{Remark}

\newcommand{\Vir}{\operatorname{Vir}}
\newcommand{\PP}{\operatorname{PP}}
\newcommand{\ord}{\operatorname{ord}}
\newcommand{\rank}{\operatorname{rank}}

\title{A pole-free torus one-point Virasoro-block recursion
at \(c=\frac12\)}
\author{}
\date{}

\begin{document}
\maketitle

\begin{abstract}
We construct a coefficientwise finite recursion for the eta-reduced torus
one-point Virasoro block directly at the Ising central charge.  Degenerate
weights and their collision classes are classified exactly.  At every grade,
generic Kac recursion terms are grouped by their common Ising weight, and
each fixed-central-charge Laurent coefficient is an exact collision moment:
the constant term of a finite sum of Laurent germs.  A valuation bound proves
that only finitely many moments occur, and a grade-decreasing shifted
recursion computes them.  Rational continuation and partial-fraction
uniqueness prove coefficientwise equality with the inverse-Shapovalov
definition to all orders in \(q\).  Local Smith and crossing-form results
identify the representation-theoretic meaning of the moments.  The first
collision gives a simple recursive residue at grade four; at grade ten its
embedding diamond produces one double pole, whose complete principal part
is evaluated explicitly.  The quotient identity specialization recovers
the three Ising characters from the corresponding BGG resolutions.
\end{abstract}

\section{Introduction}

For generic central charge, the eta-reduced torus one-point Virasoro block
obeys the residue recursion derived in
\cite[Section~3, equations~(7)--(9)]{HadaszJaskolskiSuchanek2010}.
At rational values of the Coulomb-gas parameter, distinct Kac labels have
the same degenerate weight and the individual generic-recursion summands
develop parameter poles.  Javerzat, Santachiara, and Foda analyzed the
cancellation of these non-physical poles in examples and in the
identity-insertion character problem
\cite[Sections~4--6, especially equations~(60)--(61)]%
{JaverzatSantachiaraFoda2018}.  Stocco subsequently constructed
central-charge-pole-free torus expressions through fourth order
\cite[Sections~II.B--II.C and Appendix~A]{Stocco2022}.  Ribault identifies
a recursion at rational \(\beta^2\) as an open problem and singles out
torus one-point Virasoro blocks as its simplest incarnation
\cite[Section~4.3.3]{Ribault2026}.

This paper solves that problem at \(c=1/2\) for the meromorphic Verma-module
block.  ``Pole-free'' here means free of the divergent, non-physical
parameter poles of the generic recursion; the genuine Laurent poles in the
internal weight \(\Delta\) remain.  At every grade we partition the generic
Kac summands into their exact Ising collision classes.  Each class has a
regular sum, and its fixed-\(c\) coefficients are finite constant terms of
finite Laurent sums.  The shifted recursion decreases the remaining grade,
while an explicit valuation bound limits both the pole multiplicity and the
number of Laurent coefficients required.  Theorem~\ref{thm:all-order-recursion}
proves that this construction equals the inverse-Shapovalov definition at
every power of \(q\).

The local calculations exhibit how the algebraic confluence records
Virasoro embedding geometry.  The first collision is transverse and gives
the level-four residue in
Theorem~\ref{thm:first-recursive-residue}.  The first embedding diamond
appears at level ten, creates a double pole, and is evaluated completely in
Theorem~\ref{thm:grade-ten-principal-part}.  Finally,
Theorem~\ref{thm:ising-characters} distinguishes the Verma identity
specialization from the irreducible quotient and recovers all three Ising
characters.

\section{Blocks and the reduced series}

\begin{definition}[Virasoro Verma module]
The Virasoro algebra has generators \(L_n\), \(n\in\mathbb Z\), and central
element \(c\), with
\[
 [L_m,L_n]=(m-n)L_{m+n}
 +\frac{c}{12}(m^3-m)\delta_{m+n,0}.
\]
The Verma module \(M(c,\Delta)\) is generated by a vector
\(\lvert\Delta\rangle\) such that
\[
 L_0\lvert\Delta\rangle=\Delta\lvert\Delta\rangle,\qquad
 L_n\lvert\Delta\rangle=0\quad(n>0).
\]
For a partition \(Y=(y_1\geq\cdots\geq y_\ell)\), set
\(L_{-Y}=L_{-y_1}\cdots L_{-y_\ell}\), \(|Y|=\sum_i y_i\), and
\(\ell(Y)=\ell\).
\end{definition}

\begin{definition}[Shapovalov and torus matrices]\label{def:matrices}
At level \(N\), let
\[
 G_N(c,\Delta)_{Y,Y'}
 =\langle\Delta|L_YL_{-Y'}|\Delta\rangle ,
 \qquad |Y|=|Y'|=N,
\]
where \(L_Y=L_{y_\ell}\cdots L_{y_1}\).
For a primary field \(V_d(z)\) normalized by
\(\langle\Delta|V_d(z)|\Delta\rangle=z^{-d}\), let
\[
 \rho_N(c,\Delta,d)_{Y,Y'}
 =\langle\Delta|L_YV_d(1)L_{-Y'}|\Delta\rangle .
\]
Whenever \(G_N\) is invertible, define
\[
 F_N(c,\Delta,d)
 =\sum_{|Y|=|Y'|=N}
 \bigl(G_N(c,\Delta)^{-1}\bigr)_{Y,Y'}
 \rho_N(c,\Delta,d)_{Y,Y'} .
\]
\end{definition}

\begin{definition}[Eta-reduced block]
Write
\[
 \mathcal F_\Delta^{(c)}(d\mid q)
 =q^{\Delta-c/24}\sum_{N\geq0}F_N(c,\Delta,d)q^N
\]
and define
\[
 H_\Delta^{(c)}(d\mid q)
 =\prod_{m\geq1}(1-q^m)
   \sum_{N\geq0}F_N(c,\Delta,d)q^N
 =\sum_{N\geq0}H_N(c,\Delta,d)q^N .
\]
\end{definition}

\begin{lemma}[Polynomial matrix entries]\label{lem:polynomial-entries}
For fixed \(N\), every entry of \(G_N(c,\Delta)\) belongs to
\(\mathbb Q[c,\Delta]\), and every entry of
\(\rho_N(c,\Delta,d)\) belongs to \(\mathbb Q[c,\Delta,d]\).
\end{lemma}

\begin{proof}
Repeated use of the Virasoro commutator moves every positive mode in
\(L_Y\) to the right.  Positive modes annihilate
\(\lvert\Delta\rangle\), while each \(L_0\) and the central element act by
\(\Delta\) and \(c\), respectively.  This gives the first assertion.

For the second assertion use, in addition,
\[
 [L_n,V_d(z)]
 =z^n\bigl(z\partial_z+(n+1)d\bigr)V_d(z).
\]
Move positive modes through \(V_d(z)\), then move the remaining negative
modes through \(V_d(z)\) toward the bra.  The process terminates because
each commutation removes one mode from one side of the insertion.
The remaining derivatives act on \(z^{-d}\); after setting \(z=1\), their
values are polynomials in \(d\).  The other commutators contribute only
polynomials in \(c\) and \(\Delta\).
\end{proof}

\section{Ising resonance classes}

\begin{definition}[Ising Kac weights]
For \(r,s\in\mathbb Z_{>0}\), define
\[
 \Delta_{r,s}
 =\frac{(4r-3s)^2-1}{48}.
\]
Set
\[
 \lambda_n=\frac{n^2-1}{48}\qquad(n\in\mathbb Z_{\geq0}).
\]
\end{definition}

\begin{proposition}[Exact collision classification]\label{prop:collisions}
For positive pairs \((r,s)\) and \((r',s')\),
\(\Delta_{r,s}=\Delta_{r',s'}\) if and only if one of the following holds
for some \(\ell\in\mathbb Z\):
\[
 (r',s')=(r+3\ell,s+4\ell),
\]
or
\[
 (r',s')=(3\ell-r,4\ell-s).
\]
\end{proposition}

\begin{proof}
The equality of the two weights is equivalent to
\[
 4r'-3s'=\varepsilon(4r-3s)
\]
for some \(\varepsilon\in\{1,-1\}\).
If \(\varepsilon=1\), then
\(4(r'-r)=3(s'-s)\).  Since \(\gcd(3,4)=1\), there is an
\(\ell\in\mathbb Z\) with
\((r'-r,s'-s)=(3\ell,4\ell)\).
If \(\varepsilon=-1\), then
\(4(r'+r)=3(s'+s)\), and the same coprimality argument gives
\((r'+r,s'+s)=(3\ell,4\ell)\).
The converse follows by direct substitution.
\end{proof}

\begin{corollary}[Distinct weights and local finiteness]\label{cor:local-finite}
The set of distinct positive-label degenerate weights at \(c=1/2\) is
\[
 \{\lambda_n:n\in\mathbb Z_{\geq0}\}.
\]
Every collision class contains infinitely many positive pairs.  For each
fixed level \(N\), only finitely many classes contain a pair with \(rs\leq N\).
\end{corollary}

\begin{proof}
For every positive pair, \(\Delta_{r,s}=\lambda_{|4r-3s|}\).
Conversely, given \(n\geq0\), choose a sufficiently large positive
\(r\equiv n\pmod 3\); then
\(s=(4r-n)/3\) is a positive integer and \(4r-3s=n\).
Starting with any positive pair in a class, the pairs
\((r+3\ell,s+4\ell)\), \(\ell\geq0\), are positive and remain in the
same class.  Finally, \(rs\leq N\) implies \(r\leq N\) and \(s\leq N\),
so only finitely many pairs, and hence classes, occur.
\end{proof}

\section{Coefficientwise partial fractions}

\begin{lemma}[Large internal weight]\label{lem:large-weight}
For fixed \(c,d,N\),
\[
 \lim_{\Delta\to\infty}F_N(c,\Delta,d)=p(N),
\]
where \(p(N)\) is the number of partitions of \(N\).  Consequently,
\[
 \lim_{\Delta\to\infty}H_N(c,\Delta,d)=\delta_{N,0}.
\]
\end{lemma}

\begin{proof}
For a partition \(Y\), put \(a_Y=\ell(Y)\), and rescale the PBW vector
\(L_{-Y}\lvert\Delta\rangle\) by \(\Delta^{-a_Y/2}\).
Commuting positive modes to the right shows that
\[
 \deg_\Delta G_N(Y,Y')\leq\min(a_Y,a_{Y'}).
\]
If \(a_Y=a_{Y'}\) and \(Y\neq Y'\), the degree is at most \(a_Y-1\).
For \(Y=Y'\), the leading term is
\[
 \left(\prod_{j\geq1}m_j(Y)!(2j)^{m_j(Y)}\right)
 \Delta^{a_Y},
\]
where \(m_j(Y)\) is the multiplicity of \(j\) in \(Y\).
Indeed, maximal degree requires every positive mode to be paired with a
negative mode of the same index, using the \(2jL_0\) term of
\([L_j,L_{-j}]\); any different pairing uses at least one commutator
without an \(L_0\) contribution.

Apply the same commutation argument with one \(V_d\) insertion.  The term
in which no Ward commutator is used has the same leading term as the
Shapovalov entry.  Every term using
\([L_n,V_d(z)]\) loses one possible positive--negative pairing.
It follows that, after the above rescaling, both \(G_N\) and \(\rho_N\)
converge to the same invertible diagonal matrix \(D_N\).
Therefore
\[
 \lim_{\Delta\to\infty}F_N(c,\Delta,d)
 =\operatorname{tr}(D_N^{-1}D_N)=p(N).
\]
If \(e_j\) are the coefficients of
\(\prod_{m\geq1}(1-q^m)\), then
\[
 H_N=\sum_{j=0}^N e_jF_{N-j}.
\]
The generating-function identity
\[
 \left(\sum_{j\geq0}e_jq^j\right)
 \left(\sum_{k\geq0}p(k)q^k\right)=1
\]
gives the second limit.
\end{proof}

\begin{theorem}[Finite coefficientwise pole expansion]\label{thm:partial-fractions}
At \(c=1/2\), \(H_0=1\), and for every \(N\geq1\) there are finite
integers \(m_{n,N}\) and coefficients \(a_{n,k,N}(d)\in\mathbb Q[d]\)
such that
\[
 H_N\left(\frac12,\Delta,d\right)
 =\sum_{\substack{n\geq0\\
       \exists\,r,s>0:\ rs\leq N,\ |4r-3s|=n}}
   \ \sum_{k=1}^{m_{n,N}}
   \frac{a_{n,k,N}(d)}{(\Delta-\lambda_n)^k}.
\]
In particular, setting
\[
 A_{n,k}(d;q)=\sum_{N\geq1}a_{n,k,N}(d)q^N
\]
gives a coefficientwise locally finite expansion
\[
 H_\Delta^{(1/2)}(d\mid q)
 =1+\sum_{n\geq0}\sum_{k\geq1}
 \frac{A_{n,k}(d;q)}{(\Delta-\lambda_n)^k}.
\]
\end{theorem}

\begin{proof}
By Lemma~\ref{lem:polynomial-entries},
\(F_N=\operatorname{tr}(G_N^{-1}\rho_N^{\mathsf T})\) is rational in
\(\Delta\), and its poles are contained in the zero locus of
\(\det G_N\).  The Kac determinant formula
\cite[Theorem~4.2 and equations (4.16)--(4.21)]{IoharaKoga2011} gives,
up to a nonzero factor independent of \(\Delta\),
\[
 \det G_N(c,\Delta)
 =C_N\prod_{\substack{r,s>0\\rs\leq N}}
   \bigl(\Delta-\Delta_{r,s}(c)\bigr)^{p(N-rs)}.
\]
At \(c=1/2\), Proposition~\ref{prop:collisions} identifies the distinct
zeros with the stated finite set of \(\lambda_n\)'s.
Since \(H_N\) is a finite integer linear combination of
\(F_0,\ldots,F_N\), it has the same pole containment.
Lemma~\ref{lem:large-weight} shows that \(H_N\) tends to zero at infinity
for \(N\geq1\), so its partial-fraction expansion has no polynomial part.
All matrix entries have rational coefficients at \(c=1/2\), and their
dependence on \(d\) is polynomial; hence the partial-fraction coefficients
belong to \(\mathbb Q[d]\).  Local finiteness follows from
Corollary~\ref{cor:local-finite}.
\end{proof}

\section{Local Smith exponents and the first collision}

\begin{lemma}[Pole order from local invariant factors]\label{lem:smith}
Let \(R=\mathbb C[[x]]\), and let \(G(x)\in M_m(R)\) have nonzero
determinant.  Suppose the Smith exponents of \(G\) are
\(0\leq e_1\leq\cdots\leq e_m\).  Then
\[
 \ord_x\det G=\sum_{i=1}^m e_i,\qquad
 \dim\ker G(0)=\#\{i:e_i>0\},
\]
and the pole order of \(G(x)^{-1}\) is \(e_m\).
In particular, if
\[
 \ord_x\det G=\dim\ker G(0)=k,
\]
then every positive Smith exponent equals \(1\), and \(G^{-1}\) has at
most a simple pole.
\end{lemma}

\begin{proof}
Since \(R\) is a discrete valuation ring, there are
\(U,V\in\operatorname{GL}_m(R)\) such that
\[
 UGV=\operatorname{diag}(x^{e_1},\ldots,x^{e_m}).
\]
Taking determinants gives the first identity.  Reducing modulo \(x\)
gives the second.  Inverting the displayed Smith form gives pole order
\(e_m\).  If the sum of the \(k\) positive integers \(e_i\) is \(k\),
then each of them is \(1\).
\end{proof}

\begin{proposition}[Transverse inverse residue]\label{prop:transverse-residue}
Let \(V\) be finite dimensional, let
\(G(x)\colon V\to V^*\) be a symmetric matrix over \(\mathbb C[[x]]\)
with nonzero determinant, and put \(K=\ker G(0)\).  Let
\(\iota\colon K\hookrightarrow V\) be the inclusion and define the crossing
form
\[
 \beta(u,v)=u^{\mathsf T}G'(0)v,\qquad u,v\in K.
\]
If \(\beta\) is nondegenerate, then \(G(x)^{-1}\) has only a simple pole and
\[
 \operatorname*{Res}_{x=0}G(x)^{-1}
 =\iota\,\beta^{-1}\iota^*.
\]
Moreover, \(\beta\) is nondegenerate if and only if
\[
 \ord_x\det G=\dim K.
\]
\end{proposition}

\begin{proof}
Choose a complement \(C\) of \(K\).  Since \(K\) is the radical of the
symmetric form \(G(0)\), its restriction \(A\) to \(C\) is nondegenerate.
In the decomposition \(V=K\oplus C\), one therefore has
\[
 G(x)=
 \begin{pmatrix}
  xB+O(x^2)&xD+O(x^2)\\
  xD^{\mathsf T}+O(x^2)&A+O(x)
 \end{pmatrix},
\]
where \(B\) is the matrix of \(\beta\).  The Schur complement of the
lower-right block is
\[
 S(x)=xB+O(x^2).
\]
If \(B\) is invertible, block inversion gives
\[
 G(x)^{-1}
 =\frac1x
 \begin{pmatrix}B^{-1}&0\\0&0\end{pmatrix}+O(1),
\]
which is the coordinate expression of the asserted residue.  Also
\[
 \det G(x)=\det(A+O(x))\det S(x).
\]
Every entry of \(S\) is divisible by \(x\), so its determinant is divisible
by \(x^{\dim K}\), with equality of orders precisely when
\(\det B\neq0\).  This proves the final equivalence.
\end{proof}

\begin{proposition}[Two-step inverse principal part]
\label{prop:two-step-residue}
Let
\[
 G(x)=G_0+xG_1+x^2G_2+O(x^3)
\]
be as in Proposition~\ref{prop:transverse-residue}, and put
\[
 K_1=\ker G_0,\qquad
 \beta_1(u,v)=u^{\mathsf T}G_1v,\qquad
 K_2=\ker(\beta_1\colon K_1\to K_1^*).
\]
For \(v\in K_2\), choose \(v_1\in V\) satisfying
\[
 G_0v_1=-G_1v
\]
and define
\[
 \beta_2(u,v)=u^{\mathsf T}(G_2v+G_1v_1),
 \qquad u,v\in K_2.
\]
Then \(v_1\) exists, \(\beta_2\) is independent of its choice, and
\(\beta_2\) is symmetric.  If \(\beta_2\) is nondegenerate, then all Smith
exponents are at most \(2\) and
\[
 [x^{-2}]G(x)^{-1}=\iota_2\beta_2^{-1}\iota_2^*,
\]
where \(\iota_2\colon K_2\hookrightarrow V\).  In this case,
\[
 \#\{i:e_i\geq1\}=\dim K_1,\qquad
 \#\{i:e_i\geq2\}=\dim K_2.
\]
\end{proposition}

\begin{proof}
Because \(G_0\) is symmetric,
\(\operatorname{im}G_0\) is the annihilator of \(K_1\).
For \(v\in K_2\) and \(u\in K_1\),
\[
 u^{\mathsf T}G_1v=\beta_1(u,v)=0.
\]
Thus \(G_1v\in\operatorname{im}G_0\), proving the existence of \(v_1\).
Two choices differ by an element \(k\in K_1\), and
\(u^{\mathsf T}G_1k=\beta_1(u,k)=0\) for \(u\in K_2\), proving
choice-independence.

Choose a complement to \(K_1\) as in the proof of
Proposition~\ref{prop:transverse-residue}, and first take the Schur
complement of its nondegenerate block.  This produces a symmetric matrix on
\(K_1\) of the form
\[
 S(x)=xB_1+x^2B_2+O(x^3),
\]
where \(B_1\) represents \(\beta_1\), and the restriction of \(B_2\) to
\(\ker B_1=K_2\) is the displayed form \(\beta_2\).  This also proves its
symmetry.  Now choose a complement \(E\) to \(K_2\) in \(K_1\).
The restriction of \(B_1\) to \(E\) is nondegenerate, and a second Schur
complement gives
\[
 S(x)\sim
 \begin{pmatrix}
  x^2\beta_2+O(x^3)&0\\
  0&xB_1|_E+O(x^2)
 \end{pmatrix}.
\]
If \(\beta_2\) is nondegenerate, block inversion yields the asserted
\(x^{-2}\) coefficient.  The same block form gives the two counts of Smith
exponents.
\end{proof}

\begin{corollary}[Determinant multiplicity does not determine block pole
order]\label{cor:det-overcounts}
At a collided Kac weight, the multiplicity of the corresponding factor in
\(\det G_N\) is only the sum of the positive local Smith exponents.
The inverse Shapovalov pole order is their maximum.
\end{corollary}

\begin{proof}
Apply Lemma~\ref{lem:smith} with \(x=\Delta-\lambda_n\).
\end{proof}

\begin{lemma}[Torus form on a singular submodule]
\label{lem:singular-factorization}
Let \(\chi\in M(c,h)\) be a singular vector of level \(\ell\), normalized
by a fixed PBW coefficient, and let
\[
 P_\chi(d)=\langle\chi|V_d(1)|\chi\rangle .
\]
For partitions \(Y,Y'\) of the same integer \(m\),
\[
 \langle\chi|L_YV_d(1)L_{-Y'}|\chi\rangle
 =P_\chi(d)\,
 \rho_m(c,h+\ell,d)_{Y,Y'}.
\]
\end{lemma}

\begin{proof}
The assignment
\[
 L_{-Y}|h+\ell\rangle\longmapsto L_{-Y}\chi
\]
is a homomorphism from \(M(c,h+\ell)\) to \(M(c,h)\), because \(\chi\)
is singular and has \(L_0\)-weight \(h+\ell\).  Pull the torus
three-point form on \(M(c,h)\) back along this homomorphism in both module
arguments.  Commuting one Virasoro mode at a time through \(V_d(z)\) shows
that the pulled-back form obeys the same Ward recursion as
\(\rho_m(c,h+\ell,d)\).  That recursion terminates at the highest-weight
matrix element and therefore determines every descendant matrix element
uniquely.  The pulled-back highest-weight value is \(P_\chi(d)\), whereas
the normalization
\(\langle h+\ell|V_d(1)|h+\ell\rangle=1\) gives highest-weight value
\(1\) for \(\rho_m\).  The displayed identity follows.
\end{proof}

\begin{lemma}[First crossing form on a singular submodule]
\label{lem:singular-crossing}
Fix \(c\) and let \(\chi=P|h_0\rangle\in M(c,h_0)\) be singular of level
\(\ell\), where the PBW operator \(P\in U(\Vir_-)\) is held fixed as \(h\)
varies.  Set
\[
 b_\chi=
 \left.\frac{d}{dh}
 \langle h|\omega(P)P|h\rangle\right|_{h=h_0},
\]
where \(\omega(L_n)=L_{-n}\).  For partitions \(Y,Y'\) of \(m\),
\[
 \left.\frac{d}{dh}
 \langle h|\omega(P)L_YL_{-Y'}P|h\rangle\right|_{h=h_0}
 =b_\chi\,G_m(c,h_0+\ell)_{Y,Y'}.
\]
\end{lemma}

\begin{proof}
The derivative on the left defines a bilinear form \(C\) on the PBW
space of \(M(c,h_0+\ell)\).  At \(h=h_0\), the map
\[
 j_\chi\colon M(c,h_0+\ell)\longrightarrow M(c,h_0),
 \qquad u|h_0+\ell\rangle\longmapsto u\chi ,
\]
is a module homomorphism, and its image lies in the radical of the
Shapovalov form.

Negative modes intertwine with \(j_\chi\) identically.  When a positive
mode is commuted through \(P|h\rangle\), the difference from the
intertwining identity vanishes at \(h_0\) and is therefore divisible by
\(h-h_0\).  Its first-order contribution to \(C\) pairs with
\(\operatorname{im}j_\chi\) in the specialized Shapovalov form and hence
vanishes.  Thus \(C\) is contravariant.  A contravariant form on a Verma
module is uniquely determined by its highest-weight value, because PBW
commutation recursively reduces every matrix element to that value.
Here the highest-weight value is \(b_\chi\), so
\(C=b_\chi G_m(c,h_0+\ell)\).
\end{proof}

\begin{proposition}[First Ising collision]\label{prop:first-collision}
At level \(4\), the labels \((1,2)\) and \((2,2)\) collide at
\(\Delta=1/16\).  In the PBW basis
\[
 (L_{-4},\,L_{-3}L_{-1},\,L_{-2}^2,\,
   L_{-2}L_{-1}^2,\,L_{-1}^4)\lvert\Delta\rangle,
\]
the Shapovalov determinant at \(c=1/2\) is
\[
 12\Delta^3(2\Delta-7)(2\Delta-1)^3(3\Delta-5)
 (16\Delta-21)(16\Delta-1)^3.
\]
At \(\Delta=1/16\), the local Smith exponents are
\[
 (0,0,1,1,1).
\]
Consequently, \(G_4(1/2,\Delta)^{-1}\), \(F_4(1/2,\Delta,d)\), and
\(H_4(1/2,\Delta,d)\) have at most a simple pole at \(\Delta=1/16\).
\end{proposition}

\begin{proof}
Direct Virasoro commutation gives
\begin{center}
\resizebox{\textwidth}{!}{$
G_4=
\begin{pmatrix}
8\Delta+\frac52&14\Delta&24\Delta+\frac32&36\Delta&120\Delta\\
14\Delta&12\Delta^2+14\Delta&30\Delta&40\Delta^2+20\Delta&
192\Delta^2+48\Delta\\
24\Delta+\frac32&30\Delta&32\Delta^2+36\Delta+\frac{17}{8}&
48\Delta^2+51\Delta&216\Delta^2+144\Delta\\
36\Delta&40\Delta^2+20\Delta&48\Delta^2+51\Delta&
32\Delta^3+118\Delta^2+33\Delta&
288\Delta^3+336\Delta^2+96\Delta\\
120\Delta&192\Delta^2+48\Delta&216\Delta^2+144\Delta&
288\Delta^3+336\Delta^2+96\Delta&
384\Delta^4+1152\Delta^3+1056\Delta^2+288\Delta
\end{pmatrix}
$}
\end{center}
Its determinant is the displayed polynomial.  At \(\Delta=1/16\), the
upper-left \(2\times2\) minor has determinant \(2\), while the following
three independent vectors lie in the kernel:
\[
\begin{pmatrix}-9/16\\-3/2\\1\\0\\0\end{pmatrix},\qquad
\begin{pmatrix}-27/64\\-9/8\\0\\1\\0\end{pmatrix},\qquad
\begin{pmatrix}-465/256\\-75/32\\0\\0\\1\end{pmatrix}.
\]
Thus the specialized matrix has rank \(2\) and nullity \(3\).
The determinant has order \(3\) at \(1/16\), so
Lemma~\ref{lem:smith} gives the stated Smith exponents and inverse pole
bound.  The definitions of \(F_4\) and \(H_4\) then give the same bound
for the block coefficients.
\end{proof}

\begin{proposition}[Singular data at the first collision]
\label{prop:first-singular-data}
Put \(h_0=1/16\).  The following vectors in \(M(1/2,h_0)\) are singular:
\[
 \chi_2=
 \left(L_{-1}^2-\frac34L_{-2}\right)|h_0\rangle
\]
and
\[
 \chi_4=
 \left(
 L_{-1}^4-\frac{25}{6}L_{-2}L_{-1}^2
 +\frac{49}{144}L_{-2}^2+\frac{11}{6}L_{-3}L_{-1}
 -\frac14L_{-4}
 \right)|h_0\rangle .
\]
At level \(4\),
\[
 \ker G_4(1/2,h_0)
 =\operatorname{span}
 \{L_{-2}\chi_2,L_{-1}^2\chi_2,\chi_4\}.
\]
In this ordered basis the crossing form of
Proposition~\ref{prop:transverse-residue}, with \(x=\Delta-h_0\), is
\[
 \beta_4=
 \begin{pmatrix}
 -119/8&-693/32&0\\
 -693/32&-9471/128&0\\
 0&0&13475/648
 \end{pmatrix}
 =
 \begin{pmatrix}
 -\frac74G_2(1/2,33/16)&0\\
 0&13475/648
 \end{pmatrix}.
\]
Finally, the two highest-weight torus matrix elements are
\[
 P_2(d):=\langle\chi_2|V_d(1)|\chi_2\rangle
 =\frac{d(d-1)(2d-7)(2d-1)}4
\]
and
\[
 \begin{split}
 P_4(d)&:=\langle\chi_4|V_d(1)|\chi_4\rangle\\
 &=\frac{
 d(d-1)(2d-7)(2d-1)(3d-14)(3d-5)(6d-55)(6d-1)}
 {1296}.
 \end{split}
\]
\end{proposition}

\begin{proof}
Apply \(L_1,\ldots,L_\ell\) to the displayed level-\(\ell\) vector and
use the Virasoro commutator.  Collecting coefficients in the PBW bases at
levels \(\ell-1,\ldots,0\) gives zero for \(\chi_2\) and for \(\chi_4\);
modes larger than \(\ell\) vanish by degree.  Thus both vectors are
singular.

In the level-four PBW basis of
Proposition~\ref{prop:first-collision}, the three asserted kernel vectors
are the columns
\[
 \begin{pmatrix}0\\0\\-3/4\\1\\0\end{pmatrix},\qquad
 \begin{pmatrix}-3/2\\-3/2\\0\\-3/4\\1\end{pmatrix},\qquad
 \begin{pmatrix}-1/4\\11/6\\49/144\\-25/6\\1\end{pmatrix}.
\]
Multiplication by the displayed matrix \(G_4(1/2,h_0)\) gives zero.
They are independent, and Proposition~\ref{prop:first-collision} gives
nullity three, proving the kernel identity.

Multiplying these columns on both sides of
\(\partial_\Delta G_4(1/2,\Delta)|_{\Delta=h_0}\) gives the displayed
crossing matrix.  Direct level-two commutation gives
\[
 G_2(1/2,33/16)=
 \begin{pmatrix}17/2&99/8\\99/8&1353/32\end{pmatrix},
\]
which verifies the second expression for \(\beta_4\).
Finally, repeated use of the primary-field Ward commutator in the proof of
Lemma~\ref{lem:polynomial-entries}, followed by collection and
factorization over \(\mathbb Q[d]\), gives the displayed \(P_2\) and
\(P_4\).
\end{proof}

\begin{theorem}[First fixed-\(c\) recursive principal part]
\label{thm:first-recursive-residue}
At \(c=1/2\), the level-four reduced coefficient satisfies the finite
fixed-central-charge recursion
\[
 \operatorname*{Res}_{\Delta=1/16}H_4(1/2,\Delta,d)
 =-\frac47P_2(d)H_2(1/2,33/16,d)
  +\frac{648}{13475}P_4(d).
\]
Equivalently, this residue is
\[
 \frac{d(d-1)(2d-7)(2d-1)}{40425}
 \left(
 248d^4-5752d^3+25274d^2-33630d+5775
 \right).
\]
Every quantity on the right is evaluated directly at \(c=1/2\).
\end{theorem}

\begin{proof}
At levels \(2\), \(3\), and \(4\), the radicals at
\(\Delta=h_0\) are respectively
\[
 \operatorname{span}\{\chi_2\},\qquad
 \operatorname{span}\{L_{-1}\chi_2\},\qquad
 \operatorname{span}\{L_{-2}\chi_2,L_{-1}^2\chi_2,\chi_4\}.
\]
The crossing forms on the first two spaces are
\[
 -\frac74,\qquad -\frac74G_1(1/2,33/16),
\]
and the third is the form in
Proposition~\ref{prop:first-singular-data}.  These identities follow by
differentiating the level-two, level-three, and level-four Gram matrices;
all three crossing forms are nondegenerate.  Apply
Proposition~\ref{prop:transverse-residue}, then use
Lemma~\ref{lem:singular-factorization} on each singular summand.  Taking
the trace against the corresponding torus matrix gives
\[
 \begin{aligned}
 \operatorname*{Res}_{\Delta=h_0}F_2
   &=-\frac47P_2,\\
 \operatorname*{Res}_{\Delta=h_0}F_3
   &=-\frac47P_2F_1(1/2,33/16,d),\\
 \operatorname*{Res}_{\Delta=h_0}F_4
   &=-\frac47P_2F_2(1/2,33/16,d)
     +\frac{648}{13475}P_4.
 \end{aligned}
\]
The Euler product has coefficients \(1,-1,-1,0,0\) through degree four.
Hence \(H_4=F_4-F_3-F_2\), while
\(H_2=F_2-F_1-F_0\).  Substitution gives the recursive formula.
Inserting the explicit polynomials from
Proposition~\ref{prop:first-singular-data} and the level-two coefficient
from the inverse Shapovalov matrix gives the second displayed expression.
\end{proof}

\section{The first nontransverse embedding diamond}

\begin{definition}[Generic residue data used for algebraic confluence]
\label{def:generic-residue}
For \(b\in\mathbb C^\times\), set
\[
 c(b)=1+6(b+b^{-1})^2,\qquad
 h_{r,s}(b)=\frac{(b+b^{-1})^2-(rb+sb^{-1})^2}{4}.
\]
Put \(\Lambda^2=(b+b^{-1})^2-4d\), and define
\[
 \begin{split}
 {\cal P}_{r,s}(b,d)
 =\prod_{\substack{1\leq k\leq2r-1\\k\ {\rm odd}}}
  \prod_{\substack{1\leq l\leq2s-1\\l\ {\rm odd}}}
 &\frac{\Lambda^2-(kb+lb^{-1})^2}{4}\\
 {}\times&
 \frac{\Lambda^2-(kb-lb^{-1})^2}{4},
 \end{split}
\]
and
\[
 {\cal A}_{r,s}(b)
 =\frac12
 \prod_{\substack{1-r\leq p\leq r,\ 1-s\leq q\leq s\\
 (p,q)\neq(0,0),(r,s)}}
 (pb+qb^{-1})^{-1},\qquad
 {\cal R}_{r,s}(b,d)={\cal A}_{r,s}(b){\cal P}_{r,s}(b,d).
\]
For generic \(b\), \({\cal R}_{r,s}\) is the residue coefficient in the
eta-reduced torus recursion.
\end{definition}

\begin{proof}
The formula is equations (7)--(9) of
\cite[Section~3]{HadaszJaskolskiSuchanek2010}; the product
\({\cal P}_{r,s}\) is their product of the two torus fusion polynomials,
written without choosing a square root of \(\Lambda^2\).
\end{proof}

\begin{proposition}[Grade-ten Smith profile at \(\Delta=1/16\)]
\label{prop:grade-ten-smith}
Let \(h_0=1/16\), and let
\[
 U_2=U(\Vir_-)\chi_2,\qquad U_4=U(\Vir_-)\chi_4
\]
be the two singular submodules of \(M(1/2,h_0)\) from
Proposition~\ref{prop:first-singular-data}.  Their first intersection is
generated at total level \(10\) by a singular vector \(\chi_{10}\), and
\[
 \dim\ker G_{10}(1/2,h_0)=32.
\]
The \(42\) local Smith exponents of
\(G_{10}(1/2,h_0+x)\) are
\[
 \bigl(0^{\,10},1^{\,31},2\bigr).
\]
In particular, the inverse Shapovalov matrix has a nonzero rank-one
coefficient at order \(x^{-2}\).
\end{proposition}

\begin{proof}
The weight \(h_0\) is the interior Ising weight labelled by \((1,2)\), so
it is in Class \(R^+\), Case \(1+\), in the classification of
\cite[Section~5.1.4]{IoharaKoga2011}.
The commutative embedding diagram in
\cite[Proposition~5.4 and Figure~5.4]{IoharaKoga2011} describes the two
branches and their successive intersections.  The identification of the
first Jantzen layer with their sum, hence with the maximal proper
submodule, is
\cite[Theorem~6.3(i) and equation~(6.3)]{IoharaKoga2011}.

The two highest weights are
\[
 h_0+2=\frac{33}{16}=\lambda_{10},\qquad
 h_0+4=\frac{65}{16}=\lambda_{14}.
\]
The first singular vector inside \(U_2\) has relative level \(8\), with
Kac label \((4,2)\), and the first singular vector inside \(U_4\) has
relative level \(6\), with Kac label \((1,6)\).  Both therefore have total
level \(10\) and weight
\[
 h_0+10=\frac{161}{16}=\lambda_{22}.
\]
The embeddings are injective by
\cite[the paragraph preceding Corollary~5.2]{IoharaKoga2011}.
Their two images are proportional by uniqueness of a singular vector at a
fixed level, \cite[Proposition~5.1]{IoharaKoga2011}.  Hence the
level-ten intersection is one dimensional.  There is no earlier
intersection: a vector at the first intersection grade would be singular,
whereas the Kac determinant for \(M(1/2,33/16)\), respectively
\(M(1/2,65/16)\), has no zero below relative level \(8\), respectively
\(6\).

The grade-ten radical is the grade-ten part of \(U_2+U_4\).  Since
\[
 p(8)=22,\qquad p(6)=11,
\]
and the intersection has dimension one at this grade,
\[
 \dim\ker G_{10}(1/2,h_0)=22+11-1=32.
\]
By Proposition~\ref{prop:collisions}, the only positive Kac labels with
\(\Delta_{r,s}=h_0\) and \(rs\leq10\) are \((1,2)\) and \((2,2)\).
The Kac determinant formula therefore gives
\[
 \ord_x\det G_{10}(1/2,h_0+x)
 =p(8)+p(6)=33.
\]
There are \(p(10)=42\) Smith exponents, exactly \(32\) of them are
positive, and their sum is \(33\).  Thus \(31\) positive exponents equal
\(1\), one equals \(2\), and the remaining \(10\) vanish.
Lemma~\ref{lem:smith} gives the final assertion.
\end{proof}

\begin{lemma}[Exact constants of the grade-ten confluence]
\label{lem:grade-ten-constants}
Let \(b_0=2i/\sqrt3\), so \(c(b_0)=1/2\).  Then
\[
 \begin{array}{c|cccc}
 (r,s)&(1,2)&(2,2)&(4,2)&(1,6)\\ \hline
 {\cal A}_{r,s}(b_0)&
 -4/7&648/13475&48/32035128125&-8/320945625 .
 \end{array}
\]
If a dot denotes \(d/db\) at \(b_0\), then
\[
 \frac{\dot h_{1,2}-\dot h_{2,2}}
 {\dot h_{1,2}-\dot h_{4,2}}=\frac15,\qquad
 \frac{\dot h_{2,2}-\dot h_{1,2}}
 {\dot h_{2,2}-\dot h_{1,6}}=\frac17,
\]
and
\[
 \frac15{\cal A}_{1,2}{\cal A}_{4,2}
 =\frac17{\cal A}_{2,2}{\cal A}_{1,6}
 =-\frac{192}{1121229484375}.
\]
Moreover,
\[
 \begin{split}
 {\cal P}_{4,2}(b_0,d)
 =\frac1{20736}&d(d-20)(d-13)(d-11)(d-6)(d-1)\\
 {}\times&(2d-57)(2d-35)(2d-15)(2d-7)(2d-5)(2d-1)\\
 {}\times&(3d-14)(3d-5)(6d-55)(6d-1),
 \end{split}
\]
\[
 \begin{split}
 {\cal P}_{1,6}(b_0,d)
 =\frac1{64}&d(d-20)(d-13)(d-11)(d-6)(d-1)\\
 {}\times&(2d-57)(2d-35)(2d-15)(2d-7)(2d-5)(2d-1),
 \end{split}
\]
and
\[
 P_2(d){\cal P}_{4,2}(b_0,d)
 =P_4(d){\cal P}_{1,6}(b_0,d).
\]
\end{lemma}

\begin{proof}
Substitute \(b_0=2i/\sqrt3\) in the finite products of
Definition~\ref{def:generic-residue}.  Direct differentiation gives
\[
 c(b_0)=\frac12,\qquad c'(b_0)=\frac{7i\sqrt3}{2}\neq0,
\]
so \(b-b_0\) is a valid local deformation parameter.  The Kac-weight
derivatives are
\[
 \begin{array}{c|rrrr}
 (r,s)&(1,2)&(2,2)&(4,2)&(1,6)\\ \hline
 16\dot h_{r,s}/(i\sqrt3)&9&-7&-71&105 .
 \end{array}
\]
These four values give the two slope ratios.  Multiplication of the four
displayed \({\cal A}\)-values gives the common nested residue.
Factoring the two \({\cal P}\)-products over \(\mathbb Q[d]\) gives the
last three identities.
\end{proof}

\begin{theorem}[Leading double-pole coefficient at grade ten]
\label{thm:grade-ten-leading}
Normalize all singular vectors by setting their \(L_{-1}^{\ell}\)
coefficient equal to \(1\).  Let
\(\psi_8\in M(1/2,33/16)\) and
\(\psi_6\in M(1/2,65/16)\) be the singular vectors of levels \(8\) and
\(6\), respectively.  Then
\[
 \chi_{10}=j_{\chi_2}(\psi_8)=j_{\chi_4}(\psi_6).
\]
Define
\[
 P_8^{(33/16)}(d)=\langle\psi_8|V_d(1)|\psi_8\rangle,\qquad
 P_6^{(65/16)}(d)=\langle\psi_6|V_d(1)|\psi_6\rangle .
\]
They are the explicit polynomials
\[
 P_8^{(33/16)}(d)={\cal P}_{4,2}(b_0,d),\qquad
 P_6^{(65/16)}(d)={\cal P}_{1,6}(b_0,d)
\]
listed in Lemma~\ref{lem:grade-ten-constants}.
Write
\[
 G_{10}(1/2,h_0+x)=G_0+xG_1+x^2G_2+O(x^3),
\]
choose \(v_1\) with
\[
 G_0v_1=-G_1\chi_{10},
\]
and set
\[
 \gamma_{10}
 =\chi_{10}^{\mathsf T}(G_2\chi_{10}+G_1v_1).
\]
Then it is independent of the choice of \(v_1\), and
\[
 \gamma_{10}=-\frac{1121229484375}{192}.
\]
Consequently,
\[
 \lim_{\Delta\to h_0}
 (\Delta-h_0)^2H_{10}(1/2,\Delta,d)
 =\frac{P_2(d)P_8^{(33/16)}(d)}{\gamma_{10}}
 =\frac{P_4(d)P_6^{(65/16)}(d)}{\gamma_{10}}.
\]
This is an intrinsic fixed-\(c\) expression for the first
nontransverse Laurent coefficient.
\end{theorem}

\begin{proof}
The two compositions in the first display are nonzero singular vectors of
total level \(10\).  Proposition~5.1 of \cite{IoharaKoga2011} makes them
proportional.  Their \(L_{-1}^{10}\) coefficients are both \(1\), so they
are equal.

Let \(K_1=\ker G_0=(U_2+U_4)_{10}\), and let \(\beta_1\) be its first
crossing form.  By Lemma~\ref{lem:singular-crossing}, the restriction of
\(\beta_1\) to \(U_2\) is the scalar \(b_{\chi_2}\) times the
Shapovalov form on \(M(1/2,33/16)\).  Since \(\psi_8\) is singular,
\(\beta_1(\chi_{10},U_2)=0\).  The same argument through \(U_4\) gives
\(\beta_1(\chi_{10},U_4)=0\).  Therefore
\(\chi_{10}\in K_2=\ker\beta_1\).

Proposition~\ref{prop:grade-ten-smith} shows that exactly one Smith
exponent is at least \(2\).  Hence \(K_2\) is spanned by
\(\chi_{10}\), its second crossing form is the nonzero scalar
\(\gamma_{10}\), and Proposition~\ref{prop:two-step-residue} gives
\[
 [x^{-2}]G_{10}(1/2,h_0+x)^{-1}
 =\frac{\chi_{10}\chi_{10}^{\mathsf T}}{\gamma_{10}}.
\]
Taking the trace against the specialized torus matrix yields
\[
 [x^{-2}]F_{10}(1/2,h_0+x,d)
 =\frac{\langle\chi_{10}|V_d(1)|\chi_{10}\rangle}{\gamma_{10}}.
\]
For \(N<10\), the two singular submodules have zero intersection at grade
\(N\).  Thus the radical dimension is
\[
 p(N-2)+p(N-4),
\]
with \(p(k)=0\) for \(k<0\), and this equals the Kac-determinant valuation.
Lemma~\ref{lem:smith} therefore shows that all lower-level Smith exponents
at \(h_0\) are at most one.  Multiplication by the Euler product consequently
does not alter the \(x^{-2}\) coefficient at grade \(10\).  Finally,
Lemma~\ref{lem:singular-factorization}, applied along either side of the
embedding diamond, gives
\[
 \langle\chi_{10}|V_d(1)|\chi_{10}\rangle
 =P_2(d)P_8^{(33/16)}(d)
 =P_4(d)P_6^{(65/16)}(d).
\]
The Gram matrices and the normalized singular-vector coefficients are
rational at \(c=1/2\), so the defining linear system gives
\(\gamma_{10}\in\mathbb Q\).  It remains to evaluate this intrinsic scalar.
Apply the generic recursion of Definition~\ref{def:generic-residue}
algebraically near \(b_0\).  The nested path
\((1,2)\to(4,2)\) contributes
\[
 \frac{\dot h_{1,2}-\dot h_{2,2}}
 {\dot h_{1,2}-\dot h_{4,2}}\,
 {\cal A}_{1,2}(b_0){\cal A}_{4,2}(b_0)
 P_2(d)P_8^{(33/16)}(d)
\]
to the double pole.  Lemma~\ref{lem:grade-ten-constants} evaluates its
scalar coefficient as
\(-192/1121229484375\).  The other path gives the same result.  Comparison
with the already proved intrinsic coefficient
\(P_2P_8/\gamma_{10}\) yields the displayed value of \(\gamma_{10}\).
\end{proof}

\begin{theorem}[Complete grade-ten principal part]
\label{thm:grade-ten-principal-part}
Define the finite shifted-block parts
\[
 Q_8(d)=
 \left.
 \frac{\partial}{\partial z}
 \left\{(z-\tfrac{33}{16})H_8(1/2,z,d)\right\}
 \right|_{z=33/16}
\]
and
\[
 Q_6(d)=
 \left.
 \frac{\partial}{\partial z}
 \left\{(z-\tfrac{65}{16})H_6(1/2,z,d)\right\}
 \right|_{z=65/16}.
\]
Put
\[
 \begin{split}
 \Theta_{10}(d)
 ={}&-\frac{
 d^2(d-1)^2(2d-7)^2(2d-1)^2
 (3d-14)(3d-5)(6d-55)(6d-1)}
 {552735279561465000000}\\
 &\quad\times
 \bigl(
 424640464d^8-41506307104d^7+1664502727528d^6\\
 &\qquad
 -35552058801520d^5+438148123547677d^4
 -3142854705332170d^3\\
 &\qquad
 +12488401942129875d^2-23870935242348750d
 +14566042731240000
 \bigr).
 \end{split}
\]
Then the entire principal part of the level-ten reduced coefficient at
\(\Delta=h_0=1/16\) is
\[
 \begin{split}
 \PP_{\Delta=h_0}H_{10}(1/2,\Delta,d)
 ={}&
 -\frac{192\,P_2(d)P_8^{(33/16)}(d)}
 {1121229484375\,(\Delta-h_0)^2}\\
 &+
 \frac{
 -\frac47P_2(d)Q_8(d)
 +\frac{648}{13475}P_4(d)Q_6(d)
 +\Theta_{10}(d)}
 {\Delta-h_0}.
 \end{split}
\]
Every term on the right is finite and is evaluated at \(c=1/2\).
\end{theorem}

\begin{proof}
Let \(t=b-b_0\), and abbreviate the four labels
\[
 \alpha=(1,2),\quad \beta=(2,2),\quad
 \mu=(4,2),\quad \nu=(1,6).
\]
Set
\[
 z_\alpha=h_\alpha+2,\qquad z_\beta=h_\beta+4,\qquad
 \delta_\alpha=z_\alpha-h_\mu,\qquad
 \delta_\beta=z_\beta-h_\nu.
\]
Only the \(\alpha\)- and \(\beta\)-terms of the generic recursion can
develop a pole at \(\Delta=h_0\).  Their numerators at grade ten are
\[
 C_\alpha(b)
 ={\cal R}_\alpha(b,d)
 H_8(c(b),z_\alpha(b),d),\qquad
 C_\beta(b)
 ={\cal R}_\beta(b,d)
 H_6(c(b),z_\beta(b),d).
\]
At \((b_0,z_\alpha(b_0))\), the level-eight shifted block has one
transverse pole, labelled by \(\mu\), and no other Kac divisor through
that point.  The analogous statement holds for the level-six block and
\(\nu\).  Therefore the definitions of \(Q_8,Q_6\) give
\[
 \begin{aligned}
 H_8(c(b),z_\alpha(b),d)
 &=\frac{{\cal R}_\mu(b,d)}{\delta_\alpha(b)}
   +Q_8(d)+O(t),\\
 H_6(c(b),z_\beta(b),d)
 &=\frac{{\cal R}_\nu(b,d)}{\delta_\beta(b)}
   +Q_6(d)+O(t).
 \end{aligned}
\]

For either nested pair \((\kappa,\tau)\), expansion of
\({\cal R}_\kappa{\cal R}_\tau/\delta_\kappa\) gives finite term
\[
 T_{\kappa,\tau}
 =
 \frac{\dot{\cal R}_\kappa{\cal R}_\tau
       +{\cal R}_\kappa\dot{\cal R}_\tau}
      {\dot\delta_\kappa}
 -
 \frac{{\cal R}_\kappa{\cal R}_\tau\ddot\delta_\kappa}
      {2\dot\delta_\kappa^2},
\]
where every undotted factor on the right is evaluated at \(b_0\).
The two \(t^{-1}\) terms cancel because
\[
 \frac{{\cal R}_\alpha{\cal R}_\mu}{\dot\delta_\alpha}
 +\frac{{\cal R}_\beta{\cal R}_\nu}{\dot\delta_\beta}=0;
\]
this is Lemma~\ref{lem:grade-ten-constants} together with
\(P_2P_8=P_4P_6\).

Write \(y=\Delta-h_0\).  Expanding
\[
 \frac{C_\alpha(b)}{\Delta-h_\alpha(b)}
 +\frac{C_\beta(b)}{\Delta-h_\beta(b)}
\]
at \(t=0\), the cancelled \(t^{-1}\) numerator produces the \(y^{-2}\)
coefficient
\[
 \frac{\dot h_\alpha-\dot h_\beta}{\dot\delta_\alpha}
 {\cal R}_\alpha(b_0,d){\cal R}_\mu(b_0,d)
 =
 -\frac{192}{1121229484375}P_2(d)P_8^{(33/16)}(d).
\]
The \(y^{-1}\) coefficient is
\[
 {\cal R}_\alpha(b_0,d)Q_8(d)
 +{\cal R}_\beta(b_0,d)Q_6(d)
 +T_{\alpha,\mu}+T_{\beta,\nu}.
\]
The first two residue constants are
\(-4P_2/7\) and \(648P_4/13475\).
Differentiating the finite products in
Definition~\ref{def:generic-residue} and factoring over
\(\mathbb Q[d]\) gives
\[
 T_{\alpha,\mu}+T_{\beta,\nu}=\Theta_{10}(d).
\]
Substitution proves the stated principal part.
\end{proof}

\begin{remark}[Exact replay]
The PBW commutation, Ward recursion, determinant, rank, and null-vector
checks in Propositions~\ref{prop:first-collision} and
\ref{prop:first-singular-data}, as well as every residue identity in
Theorem~\ref{thm:first-recursive-residue}, are replayed by
\path{tools/audit_low_levels.py}.  The computation uses exact rational
arithmetic.  The Kac slopes, four generic residue constants, fusion
factorization, value of \(\gamma_{10}\), and polynomial
\(\Theta_{10}\) are replayed by
\path{tools/audit_grade10_confluence.py}.  The shifted Laurent
recursion and its first two collision moments are independently replayed
by \path{tools/audit_confluent_recursion.py}.
\end{remark}

\section{All-order confluent recursion}

\begin{definition}[Shifted Laurent recursion]\label{def:shifted-laurent}
Let \(t=b-b_0\), and work in
\[
 K((t)),\qquad K=\mathbb Q(i\sqrt3,d).
\]
For a positive Kac label \(a=(r,s)\), write
\[
 \ell_a=rs,\qquad h_a(t)=h_{r,s}(b_0+t),\qquad
 R_a(t)={\cal R}_{r,s}(b_0+t,d),\qquad z_a(t)=h_a(t)+\ell_a.
\]
Define shifted Laurent coefficients recursively by
\[
 S_{a,0}(t)=1
\]
and, for \(M\geq1\),
\[
 S_{a,M}(t)
 =\sum_{\substack{e=(u,v)\in\mathbb Z_{>0}^2\\\ell_e=uv\leq M}}
 \frac{R_e(t)}
 {z_a(t)-h_e(t)}S_{e,M-\ell_e}(t).
\]
Every denominator in this formula is a nonzero rational function of
\(b\), and the recursion terminates because its second index strictly
decreases.
\end{definition}

\begin{lemma}[No generic shifted self-poles]
\label{lem:no-generic-shifted-pole}
For positive labels \(a=(r,s)\) and \(e=(u,v)\),
\[
 h_{r,s}(b)+rs-h_{u,v}(b)
\]
is not the zero rational function of \(b\).
\end{lemma}

\begin{proof}
The Kac formula gives the identity
\[
 h_{r,s}(b)+rs=h_{r,-s}(b).
\]
If this were identically equal to \(h_{u,v}(b)\), then
\[
 (rb-sb^{-1})^2=(ub+vb^{-1})^2
\]
in \(\mathbb C(b)\).  Since this is an integral domain,
\[
 rb-sb^{-1}=\pm(ub+vb^{-1}).
\]
Comparison of the coefficients of \(b\) and \(b^{-1}\) contradicts
\(r,s,u,v>0\).  Thus every shifted denominator in
Definition~\ref{def:shifted-laurent} has a well-defined Laurent expansion.
\end{proof}

\begin{definition}[Collision moments]\label{def:collision-moments}
For \(N\geq1\), set
\[
 {\cal C}_{n,N}
 =\{a=(r,s)\in\mathbb Z_{>0}^2:
       \ell_a\leq N,\ |4r-3s|=n\}.
\]
For \(a\in{\cal C}_{n,N}\), define
\[
 B_{a,N}(t)=R_a(t)S_{a,N-\ell_a}(t),\qquad
 \delta_a(t)=h_a(t)-\lambda_n.
\]
If \(\delta_a\not\equiv0\) and \(B_{a,N}\not\equiv0\), put
\[
 e_a=\ord_t\delta_a,\qquad
 v_{a,N}=\max\{0,-\ord_tB_{a,N}\}.
\]
If \(\delta_a\not\equiv0\) and \(B_{a,N}\equiv0\), put
\(e_a=\ord_t\delta_a\) and \(v_{a,N}=0\).  If
\(\delta_a\equiv0\), set \(e_a=\infty\) and \(v_{a,N}=0\).
Define the finite structural pole bound
\[
 M_{n,N}=
 \max_{a\in{\cal C}_{n,N}}
 \begin{cases}
 1+\lfloor v_{a,N}/e_a\rfloor,&e_a<\infty,\\
 1,&e_a=\infty,
 \end{cases}
\]
and, for \(1\leq k\leq M_{n,N}\), define
\[
 A_{n,k,N}(d)
 =\operatorname{CT}_{t}
 \sum_{a\in{\cal C}_{n,N}}
 B_{a,N}(t)\delta_a(t)^{k-1}.
\]
Here \(\operatorname{CT}_t\) denotes the coefficient of \(t^0\) in a
Laurent series.  Empty collision classes contribute zero.  The actual
pole multiplicity at grade \(N\) is
\[
 m_{n,N}=\max\bigl(\{k:A_{n,k,N}\neq0\}\cup\{0\}\bigr),
\]
and satisfies \(m_{n,N}\leq M_{n,N}\).
\end{definition}

\begin{theorem}[Pole-free Ising recursion and all-order equality]
\label{thm:all-order-recursion}
For every \(N\geq1\),
\[
 H_N(1/2,\Delta,d)
 =\sum_{\substack{n\geq0\\{\cal C}_{n,N}\neq\varnothing}}
   \sum_{k=1}^{M_{n,N}}
   \frac{A_{n,k,N}(d)}{(\Delta-\lambda_n)^k}.
\]
The sums and every constant-term calculation are finite,
\(A_{n,k,N}(d)\in\mathbb Q[d]\), and the recursion in
Definitions~\ref{def:shifted-laurent}--\ref{def:collision-moments}
terminates after grades strictly below \(N\).  Consequently, with
\[
 A_{n,k}(d;q)=\sum_{N\geq1}A_{n,k,N}(d)q^N,
\]
one has the coefficientwise locally finite identity
\[
 H_\Delta^{(1/2)}(d\mid q)
 =1+\sum_{n\geq0}\sum_{k\geq1}
 \frac{A_{n,k}(d;q)}{(\Delta-\lambda_n)^k}.
\]
This series agrees coefficient by coefficient with
Definition~\ref{def:matrices}, hence with the inverse-Shapovalov
definition of the torus one-point block.
\end{theorem}

\begin{proof}
For generic \(b\), the eta-reduced recursion in
Definition~\ref{def:generic-residue} gives
\[
 H_N(c(b),\Delta,d)
 =\sum_{\ell_a\leq N}
 \frac{R_a(t)}
 {\Delta-h_a(t)}
 H_{N-\ell_a}(c(b),z_a(t),d).
\]
Induction on \(M\) in the same recursion gives
\[
 S_{a,M}(t)=H_M(c(b),z_a(t),d)
\]
for generic \(b\).  Thus
\[
 H_N(c(b),\Delta,d)
 =\sum_{\ell_a\leq N}
 \frac{B_{a,N}(t)}{\Delta-h_a(t)}.
\tag{6.1}\label{eq:generic-grouped}
\]
All sums are finite.

Fix \(N\), and choose \(\Delta\) outside the finite set of Ising Kac
weights occurring at grade \(N\).  By
Lemma~\ref{lem:polynomial-entries}, the inverse-Shapovalov expression is
analytic in \(t\) near zero: its specialized Gram determinant is nonzero
at this \(\Delta\).  Equation~\eqref{eq:generic-grouped}, initially valid
for generic \(b\), is an identity of rational functions and therefore
holds on the punctured formal neighborhood of \(b_0\).

Group its right-hand side according to the distinct limits
\(\lambda_n\):
\[
 {\cal G}_{n,N}(t,\Delta)
 =\sum_{a\in{\cal C}_{n,N}}
 \frac{B_{a,N}(t)}{\Delta-h_a(t)}.
\]
Each coefficient of a negative power of \(t\) in
\({\cal G}_{n,N}\) is a rational function of \(\Delta\), tends to zero
as \(\Delta\to\infty\), and has its poles only at
\(\Delta=\lambda_n\).  The sum of these negative Laurent coefficients
over \(n\) vanishes because the left-hand side of
\eqref{eq:generic-grouped} is analytic in \(t\).
Partial-fraction uniqueness for the distinct \(\lambda_n\)'s forces the
negative coefficient of each \({\cal G}_{n,N}\) to vanish separately.
Hence every collision class has a finite algebraic confluent value.

Put \(x=\Delta-\lambda_n\).  In the formal expansion at \(x=\infty\),
\[
 \frac1{\Delta-h_a(t)}
 =\frac1{x-\delta_a(t)}
 =\sum_{j\geq0}\frac{\delta_a(t)^j}{x^{j+1}}.
\]
Taking \(\operatorname{CT}_t\) after summing the collision class shows
that the coefficient of \(x^{-k}\) is exactly
\[
 \operatorname{CT}_t
 \sum_{a\in{\cal C}_{n,N}}
 B_{a,N}(t)\delta_a(t)^{k-1}
 =A_{n,k,N}(d).
\]
If \(\delta_a\not\equiv0\), then the summand has positive \(t\)-valuation
whenever \(k>1+\lfloor v_{a,N}/e_a\rfloor\), so its constant term
vanishes.  If \(\delta_a\equiv0\), it vanishes identically for \(k>1\).
This proves the bound \(M_{n,N}\), the finiteness of the pole sum, and
the finite-jet nature of every constant-term calculation.  The confluent
value of each class tends to zero as \(\Delta\to\infty\) and has no finite
pole except \(\lambda_n\); therefore these finitely many moments are its
complete partial-fraction expansion, with no polynomial remainder.

Every \(R_a(t)\) is polynomial in \(d\) with coefficients rational in
\(t\), and all shifted denominators are independent of \(d\).
Therefore every collision moment is polynomial in \(d\).
Equation~\eqref{eq:generic-grouped} and the analyticity argument identify
the resulting rational function with the specialized inverse-Shapovalov
coefficient.  The latter has rational coefficients at \(c=1/2\);
uniqueness of its partial fractions improves
\(A_{n,k,N}\in K[d]\) to \(A_{n,k,N}\in\mathbb Q[d]\).
Finally, local finiteness in \(q\) is
Corollary~\ref{cor:local-finite}, completing the all-order identity.
\end{proof}

\begin{corollary}[Recovery of the first two confluent formulas]
\label{cor:recover-low-confluence}
At \(N=4\), Theorem~\ref{thm:all-order-recursion} gives
Theorem~\ref{thm:first-recursive-residue}.  At \(N=10\), its
\(\lambda_2=1/16\) collision moments give
Theorem~\ref{thm:grade-ten-principal-part}.
\end{corollary}

\begin{proof}
At \(N=4\), the class \({\cal C}_{2,4}\) consists of
\((1,2)\) and \((2,2)\), and both \(B\)-terms are regular at \(t=0\).
The first collision moment is
\[
 -\frac47P_2(d)H_2(1/2,33/16,d)
 +\frac{648}{13475}P_4(d),
\]
while all higher moments vanish.

At \(N=10\), the same first-label class has
\[
 B_{(1,2),10}
 ={\cal R}_{1,2}H_8(c(b),h_{1,2}+2,d),\qquad
 B_{(2,2),10}
 ={\cal R}_{2,2}H_6(c(b),h_{2,2}+4,d).
\]
Their zeroth and first collision moments are exactly the two Laurent
coefficients computed in the proof of
Theorem~\ref{thm:grade-ten-principal-part}; all higher moments vanish by
the grade-ten Smith profile.
\end{proof}

\section{Irreducible Ising specialization}

\begin{proposition}[Identity insertion in the Verma recursion]
\label{prop:identity-verma}
For \(d=0\),
\[
 H_\Delta^{(1/2)}(0\mid q)=1
\]
as a rational function of \(\Delta\).  Thus specialization of the
meromorphic Verma block at a degenerate weight is not the irreducible
Ising trace; the latter requires taking the quotient by the radical.
\end{proposition}

\begin{proof}
In Definition~\ref{def:generic-residue}, the factor with \(k=l=1\)
in \({\cal P}_{r,s}(b,0)\) is
\[
 \frac{(b+b^{-1})^2-(b+b^{-1})^2}{4}=0.
\]
Hence every generic residue \(R_a(t)\) vanishes identically at \(d=0\).
Definitions~\ref{def:shifted-laurent} and
\ref{def:collision-moments} then give
\(A_{n,k,N}(0)=0\) for every \(N\geq1\).
Equivalently, in the inverse-Shapovalov definition the identity insertion
has \(\rho_N=G_N\), so \(F_N=p(N)\) and eta reduction gives \(H=1\).
\end{proof}

\begin{definition}[Eta-reduced irreducible identity block]
\label{def:irreducible-identity}
For \(h\in\{0,1/2,1/16\}\), let \(L(1/2,h)\) be the irreducible quotient
of \(M(1/2,h)\), and define
\[
 \widehat H_h(q)
 =\prod_{m\geq1}(1-q^m)
  \sum_{N\geq0}\dim L(1/2,h)_{h+N}\,q^N.
\]
\end{definition}

\begin{theorem}[The three Ising characters]
\label{thm:ising-characters}
The quotient identity blocks are
\[
 \widehat H_0(q)
 =\sum_{j\in\mathbb Z}
 \left(
 q^{((24j+1)^2-1)/48}
 -q^{((24j+7)^2-1)/48}
 \right),
\]
\[
 \widehat H_{1/2}(q)
 =\sum_{j\in\mathbb Z}
 \left(
 q^{((24j+5)^2-25)/48}
 -q^{((24j+11)^2-25)/48}
 \right),
\]
and
\[
 \widehat H_{1/16}(q)
 =\sum_{j\in\mathbb Z}
 \left(
 q^{((24j-2)^2-4)/48}
 -q^{((24j+10)^2-4)/48}
 \right).
\]
Consequently,
\[
 \chi_h(q)
 =q^{h-1/48}
  \frac{\widehat H_h(q)}
  {\prod_{m\geq1}(1-q^m)}
\]
are the vacuum, energy, and spin Ising characters.  Their eta-reduced
beginnings are
\[
 \begin{aligned}
 \widehat H_0&=1-q-q^6+q^{11}+q^{13}-q^{20}+\cdots,\\
 \widehat H_{1/2}&=1-q^2-q^3+q^7+q^{17}-q^{25}-q^{28}+\cdots,\\
 \widehat H_{1/16}&=1-q^2-q^4+q^{10}+q^{14}-q^{24}-q^{30}+\cdots.
 \end{aligned}
\]
\end{theorem}

\begin{proof}
The Ising model is the minimal series \(c_{3,4}=1/2\).
The Class \(R^+\), Case \(1+\), BGG resolution in
\cite[Theorem~6.9]{IoharaKoga2011} is an exact complex of Verma modules.
Taking its Euler--Poincaré character gives, for an interior label
\((r,s)\),
\[
 \chi_{r,s}(q)
 =\frac1{\eta(q)}
 \sum_{j\in\mathbb Z}
 \left(
 q^{(24j+4r-3s)^2/48}
 -q^{(24j+4r+3s)^2/48}
 \right).
\]
Use \((r,s)=(1,1),(2,1),(1,2)\), respectively.  Since
\[
 h_{r,s}=\frac{(4r-3s)^2-1}{48},\qquad
 \eta(q)=q^{1/24}\prod_{m\geq1}(1-q^m),
\]
factoring out \(q^{h_{r,s}-1/48}\) gives exactly the three displayed
\(\widehat H\)-series.  For any fixed power of \(q\), only finitely many
integers \(j\) contribute, so the Euler--Poincaré sums are
coefficientwise finite.  Direct evaluation of the exponents through
degree \(30\) gives the final three lines.
\end{proof}

\begin{remark}[Character replay]
The three BGG numerators and the resulting quotient dimensions through
grade \(15\) are replayed by
\path{tools/audit_ising_characters.py} using exact integer arithmetic.
\end{remark}

\begin{thebibliography}{9}

\bibitem{IoharaKoga2011}
K.~Iohara and Y.~Koga,
\emph{Representation Theory of the Virasoro Algebra},
Springer Monographs in Mathematics, Springer (2011),
Theorem~4.2, Propositions~5.1 and~5.4, Figure~5.4, and
Theorems~6.3 and~6.9.

\bibitem{HadaszJaskolskiSuchanek2010}
L.~Hadasz, Z.~Jask\'olski and P.~Suchanek,
\emph{Recursive representation of the torus 1-point conformal block},
JHEP \textbf{01} (2010) 063,
\href{https://arxiv.org/abs/0911.2353}{arXiv:0911.2353}.

\bibitem{JaverzatSantachiaraFoda2018}
N.~Javerzat, R.~Santachiara and O.~Foda,
\emph{Notes on the solutions of Zamolodchikov-type recursion relations in
Virasoro minimal models},
JHEP \textbf{08} (2018) 183,
\href{https://arxiv.org/abs/1806.02790}{arXiv:1806.02790}.

\bibitem{Stocco2022}
D.~Stocco,
\emph{The torus one-point block of 2d CFT and null vectors in
\(\widehat{\mathfrak{sl}}_2\)},
\href{https://arxiv.org/abs/2209.08653}{arXiv:2209.08653} (2022).

\bibitem{Ribault2026}
S.~Ribault,
\emph{Exactly solvable conformal field theories},
\href{https://arxiv.org/abs/2411.17262v4}{arXiv:2411.17262v4} (2026).

\end{thebibliography}

\end{document}
